\chapter{Esercizi base Linux}
\label{chapter:Esercizi-base-Linux}
Nel capitolo si approfondiscono attraverso esercizi i comandi visti nei vari capitoli.





\begin{tbl}{lccc}{dddddddd}{dfsdfsd}
	\textbf{Type} & $\mathbf{N_C}$ & \textbf{$\mathbf{F_S}$  (Hz)} & \textbf{$\mathbf{N_B}$  (Bits)} \\
	\midrule
	EMG & 1 & 3200 & 24 \\
	EMG & 2 & 1600 & 24  \\
	EMG & 3 & 800 & 24 \\
	EMG & 3 & 200 & 24 \\
	ACC & 3 & 104 & 16 \\
	ACC + GYRO & 6 & 104 & 16 \\
	ACC + GYRO & 6 & 208 & 16 \\
	ACC + GYRO & 6 & 416 & 16 \\
\end{tbl}




\cite{example-website}














































% l'deale sarebbe cercare di strutturare ogni sezione in modo standard tuttavia non so se è fattibile in quanto comunque ogni comando permette di effettuare azioni abbastanza diverse sul sistemam, magari potre adottare questo approccio in modo implicito ovvero riportando sempre nello stesso ordine le informazioni sul comando e poi magari creare delle sottosezioni per aspetti particolari.

% sicuramente il comando ls, poi magari il find e il comando man e info, 

% comando mkdir, touch  echo cat mv rm e cp

% comandi per i programmi di compressione gzip bzip2 e xy e il comando tar per gestione degli archivi, eventualmente appunti sul programma zip che non ho minimamente scritto nel generale

% comandi cut less more head tail e sort, non riportati negli appunti ma nel file che leggo sono citate quando si parla di reindirizzamento di i/o.

% il comando grep, già abbastanza approfondito

% La parte di shell scripting è stata iniziata, non è stata conclusa perchè non molto interessante a fini pratici. comando chmod per concedere i privilegi ai file, non visto ma solo citato




% #######################utili

% utile per evitare che latex interpreti \verb|caratteri da non interpretare|

%  per riferirsi ad un listing \lstlistingname~\ref{base1}
%  per scrivere un listing indicizzato \begin{lstlisting}[style=shell, mathescape=true, caption=Contenuto della directory \ti{doc}, label=base2] \end{lstlisting}